\title{Miss Floors and Storage Sizing for Energy-Harvesting Sensing Nodes with Events of Random Duration}
\author{Dat~Lam~Quoc%
\thanks{Dat Lam Quoc is with the FPT School of Business and Technology, FPT University, Ho Chi Minh City, Vietnam (e-mail: lamquocdat@gmail.com; ORCID 0009-0004-5432-9343).}}
\maketitle

\begin{abstract}
A node of a low-power sensor network that senses and reports events on harvested energy must be energy-aware: how much storage and what duty cycle keep event detection within a target miss probability? Such nodes miss events between activations and during energy outages, and every activation pays for sensing and reporting. From a measured harvest trace we bound both losses. For finite storage we prove an energy-induced miss floor for events with duration and a sandwich for the smallest capacity that meets a miss target, for event-independent schedules that never enter outage. A pre-registered train/test evaluation uses \nCellsSS{} site--season cells and \nEv{} events. No floor check is violated; \nSubstPass{} of \nSubst{} substantive predictions pass, \nSubstFail{} fail in magnitude, and the \nExpl{} exploratory predictions fail. Periodic operation needs \AiiTsInvLo--\AiiTsInvHi{} times the storage of the taut-string schedule, an upper bound on the least capacity. The proved lower bound is within a median factor of \AiiMedLBTS{} of it (largest ratio \AiiMaxLBTSreg{} registered, \AiiMaxLBTS{} exact). Average-day sizing underestimates storage by a median factor of \avgHCMC{} (Ho Chi Minh City, \nCellHCMC{} cells) to \avgIn{} (indoor, \nCellIn{} setups). Seasonal storage need varied up to \AivbHi{}-fold across six years at one site.
\end{abstract}

\begin{IEEEkeywords}
Energy harvesting, energy storage sizing, harvest-aware duty cycling, low-power sensor node, event detection.
\end{IEEEkeywords}

\section{Introduction}\label{sec:intro}
\IEEEPARstart{A}{} node of a low-power sensor network that senses and reports events on harvested energy has to be energy-aware before it is deployed. Such a node cannot afford to sense and report in every slot, and it stops completely when its energy storage runs empty~\cite{yue2026iotj,vanleemput2023iotj}. Both effects cost missed events. Both are fixed at design time by two provisioning choices: the duty cycle and the size of the storage element (supercapacitor or battery). Excess storage and panel area are a cost, so the design question is how much storage and what duty cycle keep event detection within a target miss probability.

Green communications and networking already asks this question for other service metrics. Solar-powered base stations are dimensioned for a target energy-outage probability~\cite{naderi2025tgcn}, and panels and batteries of radio access networks are sized jointly with resource allocation~\cite{chen2026tgcn}. Batteryless IoT devices are analysed for energy self-sustainability and uninterrupted operating time~\cite{fan2025tgcn}. Harvesting sensors with finite batteries perform distributed detection over fading channels~\cite{ardeshiri2023tgcn}, and transmit and sleep schedules are learned for multi-hop harvesting networks~\cite{dutta2025tgcn}. In a low-power sensor network the service a node delivers is often whether an event was caught and reported. This paper gives the sizing result for that metric for one reporting node. The energy of a report enters the node model (Section~\ref{sec:model}); network-level effects stay outside the bounds (Section~\ref{sec:threats}).

Design rules for sensing nodes come from two separate traditions. Energy-neutral operation~\cite{kansal2007,buchli2014sensys,ahn2023iotj} and real-time scheduling with energy curves~\cite{moser2007rts} answer \emph{when storage never empties}. Event capture with rechargeable sensors~\cite{kar2006ton,jaggi2009wn,ren2014tpds,yau2010tosn} answers \emph{how much a duty cycle misses} under stochastic recharge. A node designer with a measured harvest trace and an event-duration law needs both at once. How small can the miss probability be at a given storage capacity, and how much storage does a miss target require?

We answer both with bounds that need only a harvest trace, the node's power figures and a sample of event durations. The starting point is the miss floor of our earlier article~\cite[Thms.~4--5]{deep3_jiot2026}, which bounds event-independent schedules without storage or harvest; we restate that floor and build on it. Table~\ref{tab:new} lists what is restated and what is new. The contributions are:
\begin{enumerate}
\item \emph{When the miss floor is tight.} Evenly spaced activations attain the floor for every duration law, and only they do (Section~\ref{sec:floor}).
\item \emph{An energy floor and a sizing sandwich.} For finite storage and schedules that never enter outage, we give a miss floor for events with duration and brackets for the smallest capacity $C_{\min}(\eps)$ that meets a miss target $\eps$ (Section~\ref{sec:energy}).
\item \emph{A sizing procedure} for deployments (Section~\ref{sec:procedure}). On the measured trace it brackets what every event-independent controller that never enters outage can achieve.
\item \emph{A pre-registered trace-driven evaluation and an open artifact.} \nSubstPass{} of \nSubst{} substantive predictions frozen before the held-out data pass and \nSubstFail{} fail in magnitude; \nCons{} further predictions are consistency checks of the code. The data correction, the exact sizing brackets and the post-hoc sensitivities are amendments frozen before their runs.
\end{enumerate}
Table~\ref{tab:rq} maps the research questions to the predictions and results. The node model is dimensionless: energies are in units of the idle-slot energy $e_{\mathrm{idle}}$, and $r$ is the energy of an activation, sensing plus report, over $e_{\mathrm{idle}}$. The results therefore apply to any platform once $r$ and the leakage are known. Absolute numbers use one microcontroller camera node with published idle and active measurements~\cite{moosmann2025tinyissimo}, which we call the reference node.

\emph{Research gap.} Energy-neutral conditions~\cite{kansal2007,moser2007rts} and operating-time guarantees~\cite{fan2025tgcn} for harvesting nodes exist. Loss or outage at finite storage has been analysed for point events~\cite{zhang2013twc,ren2015p2p}, Markov recharge~\cite{deng2022arxiv,tunc2017peva} and a fixed energy cost per use~\cite{ballestrin2026cl}. We know of no storage size that meets an event-miss target for events with duration on a measured harvest trace. We also know of no floor that holds for every event-independent schedule that never enters outage.

\begin{table}[!t]\centering
\caption{Research questions, hypotheses, results and sections}\label{tab:rq}
\footnotesize\setlength{\tabcolsep}{3pt}
\begin{tabular}{@{}p{0.46\columnwidth}lll@{}}\toprule
Research question & Hyp. & Result & Sec.\\\midrule
RQ1: least miss at a given rate; which schedules attain it? & H1 & Thm.~\ref{thm:floor} & \ref{sec:floor}, \ref{sec:holds}\\
RQ2: miss under finite storage; storage for a miss target? & H1, H4 & Thms.~\ref{thm:energy},~\ref{thm:sizing} & \ref{sec:energy}, \ref{sec:spread}\\
RQ3: how far are common sizing rules from the need? & H2, H7 & Table~\ref{tab:sizing} & \ref{sec:sizing}, \ref{sec:procedure}\\
RQ4: do harvest-following or throttling help? & H3, H5 & Cor.~\ref{cor:jensen}, Prop.~\ref{prop:blackout} & \ref{sec:spread}, \ref{sec:throttle}\\\bottomrule
\end{tabular}
\end{table}

\section{Related Work}\label{sec:related}
% related_work.tex (manuscript v3, P0R2): hand-written prose, no measured numbers. Statements about cited works
% are limited to what each source states (P0R/A8_TABLE_40.md, P3R/R16_READ_40.md, lit/). Table S13 of the supplement
% holds the positioning table. The paragraph on our earlier article moved here from the procedure section (S6 m2).
Table~S13 of the supplement positions the paper against the closest works; Table~\ref{tab:new} separates it from our earlier article.

\emph{Sizing targets continuity or a point-event loss.} The storage condition for a node that never empties its storage is classical: Moser et al.~\cite{moser2007rts} give it for energy curves, and Kansal et al.~\cite{kansal2007} give the energy-neutral sizing condition for a $(\rho,\sigma)$ harvest envelope. Prediction-based duty cycling~\cite{buchli2014sensys} and learned energy management of batteryless nodes~\cite{ahn2023iotj} address operation rather than the event-miss trade-off. Design practice sizes harvester and storage for energy balance~\cite{kuaban2025iwcmc}, and an IoT feasibility model asks whether a use case can be powered at all~\cite{vanleemput2023iotj}. Closest to our objective, Zhang et al.~\cite{zhang2013twc} design a harvesting sensor node for a target event-loss probability, Deng~\cite{deng2022arxiv} and Zareei et al.~\cite{zareei2017pimrc} give a closed-form battery size for coupled data and energy queues, and Tunc and Akar~\cite{tunc2017peva} bound the first-outage probability of a given battery and tune the sensing rate to it; these works model point events or packets in a queue under Markov recharge. At larger scale, solar-powered base stations are dimensioned for an energy-outage probability~\cite{naderi2025tgcn}, and panels and batteries of cell-free access networks are sized jointly with resource allocation to minimise cost~\cite{chen2026tgcn}. Those designs address a base station or a network and an outage or cost objective, whereas the object here is one sensing node with an event-miss target.

\emph{Event-capture analysis assumes stochastic recharge.} Kar et al.~\cite{kar2006ton} bound network-wide capture with the concavity argument of Corollary~\ref{cor:jensen}. Jaggi et al.~\cite{jaggi2009wn} and Ren et al.~\cite{ren2014tpds,ren2015p2p} derive activation policies for correlated or renewal events; the memoryless policy of~\cite{jaggi2009wn} is asymptotically optimal as the energy bucket grows. Ren and Cheng~\cite{ren2015p2p} already maximise capture under a finite capacitor and an activation cost for Poisson events; Lemma~\ref{lem:greedy} is that activation budget in fluid form, and Theorem~\ref{thm:energy} adds the duration term. Quality of monitoring~\cite{yau2010tosn} and patrolling~\cite{lin2013opre} supply the gap decomposition without energy, and learned wake-up policies for IoT event reporting model the harvest as random~\cite{ruizguirola2026iotj}. Energy self-sustainability of batteryless IoT devices has a martingale characterisation with the expected uninterrupted operating time~\cite{fan2025tgcn}; that analysis has no event process, and ours has a measured trace in place of a stochastic harvest model. Donor placement keeps energy thresholds across batteryless networks~\cite{hammoud2026tgcn}, with node position as the design variable instead of storage or duty cycle. Closed forms with finite storage exist for point arrivals: energy-outage probability with Bernoulli energy arrivals~\cite{ballestrin2026cl}, blocking and depletion of a finite fluid queue~\cite{deepak2026mcap}, and depletion bounds under storage constraints~\cite{yan2024iotj}.

\emph{Offline-optimal schedules optimise throughput.} Offline-optimal transmission with a finite battery~\cite{tutuncuoglu2012twc,yang2012tcom} maximises throughput rather than event capture, and IoT work bounds or minimises age of information under stochastic harvest~\cite{lei2025iotj,botirov2026iotj,hatami2025tgcn}. We use the energy-tunnel construction for $C_{\mathrm{TS}}$ and prove the variant we need (Lemma~\ref{lem:taut}).

\emph{Intermittent IoT systems are evaluated by measurement.} Batteryless cameras~\cite{desai2022camaroptera} show the hardware side. REACT~\cite{williams2024react} argues that a fixed buffer chosen at design time fits poorly, on the same EnHANTs traces we use, and self-sustainability studies use measured indoor light~\cite{mayer2024iotj}. Before a node is built, our bounds apply on a measured trace to every event-independent design that never enters outage, and the traces check them.

\emph{Relation to our earlier article.} Theorem~\ref{thm:floor} is restated from our earlier article on frame skipping~\cite[Thms.~4--5]{deep3_jiot2026}, and the event corpus and the affine energy model follow it. Storage, harvest and sizing are absent from that article; every result of Section~\ref{sec:energy} and the trace evaluation are new here (Table~\ref{tab:new}). It used a build of the same corpus with \VoneNev{} events in which LASIESTA frames were mis-ordered; this paper uses the time-ordered corpus of \nEv{} events (Section~\ref{sec:setup}; supplement, Section~S1).

\begin{table}[!t]\centering
\caption{Restated from our earlier article~\cite{deep3_jiot2026} and new here}\label{tab:new}
\footnotesize\setlength{\tabcolsep}{3pt}
\begin{tabular}{@{}>{\raggedright\arraybackslash}p{0.41\columnwidth}>{\raggedright\arraybackslash}p{0.36\columnwidth}>{\raggedright\arraybackslash}p{0.18\columnwidth}@{}}\toprule
Result & Status & In~\cite{deep3_jiot2026}\\\midrule
Lemma~\ref{lem:gap} (gap decomposition) & Restated & classical\\
Theorem~\ref{thm:floor} (floor $\eps^*(a)$) & Restated & Thms.~4--5\\
Lemma~\ref{lem:tight} (equality iff gaps $\{g,g+1\}$) & New & periodic case\\
Corollary~\ref{cor:jensen} (spread beats burst) & Quantified & qualitative\\
Lemma~\ref{lem:moser} (storage never empties) & Restated & --\\
Lemma~\ref{lem:greedy} (activation budget) & Known for Poisson events~\cite{ren2015p2p}$^{a}$; fluid, idle-feasible form here & --\\
Lemma~\ref{lem:taut} (taut string) & Variant; proof in the supplement & --\\
Theorem~\ref{thm:energy} (energy floor) & New for events with duration & --\\
Proposition~\ref{prop:blackout} (blackout) & New term & --\\
Theorem~\ref{thm:sizing} (sizing sandwich) & New & --\\
Remark~\ref{rem:blk} (block objective) & Conditional; not checked & --\\
Proposition~\ref{prop:gate} (gate boundary) & Elementary & other quantity\\
Evaluation with harvest and storage & New & video only\\\bottomrule
\multicolumn{3}{@{}l@{}}{$^{a}$ Read from the abstract only; see Section~\ref{sec:threats}, item (xii).}
\end{tabular}
\end{table}


\section{System Model}\label{sec:model}
\emph{Slots and energy.} Time is divided into slots on a cyclic horizon of $T$ slots; one slot = one frame period of 0.1\,s (10\,fps). In slot $t$ the node draws $e_t=1+\lt+x_t(r-1)$, where $x_t\in\{0,1\}$ indicates an activation and $\lt\ge0$ is a constant leakage; energies are in units of the idle-slot energy $e_{\mathrm{idle}}$. An activation senses, runs the detector and sends one report, so $r=(E_{\mathrm{det}}+E_{\mathrm{rep}})/e_{\mathrm{idle}}$, where $E_{\mathrm{det}}$ is the sensing energy of the slot and $E_{\mathrm{rep}}$ the energy of one report. We take the worst case of a report on every activation. The schedule then stays event-independent, and the storage need is an upper bound for any rule that reports less often. Every result below treats $r$ as a parameter; the theorems are unchanged. The harvest is $h_t\ge0$, the capacity is $C$, and the storage evolves as $B_t=\min(C,B_{t-1}+h_t-e_t)$. If a slot cannot be paid, the node is in \emph{outage}: it skips activations, draws only leakage, and restarts when $B\ge\theta C$ ($\theta=0.05$). A schedule \emph{never enters outage} if this never happens.

\emph{Policies.} A schedule $x\in\{0,1\}^T$ has \emph{realised} mean activation $a(x)=T^{-1}\sum_tx_t$, counted over all slots including outages. It is \emph{event-independent} (oblivious) if it does not depend on the event process. It may depend on time, on the harvest trace (even offline) and on its own storage. Because the report energy is charged on every activation, the storage trajectory does not depend on the events. A \emph{gate} pays a cue cost $c$ per slot and activates according to frame content.

\emph{Events.} An event has onset $s$ and duration $D\in\{1,2,\dots\}$ slots. It is caught if some activation falls in $\{s,\dots,s+D-1\}$; $\mathrm{miss}(x)$ is the probability that it is not. The detector is assumed to detect an event whenever it is invoked during it, so the event-miss probability measures scheduling only.

\emph{Assumption A1.} Onsets are uniform on the horizon, $D\sim F_D$ with $D\le T$, and $(s,D,x)$ are independent.

\emph{Notation.} $\phi(g)=\E[\min(g,D)]$ and $\Psif(a)=\E[\min(1,aD)]$ are concave and non-decreasing. $\eps^*(a)=1-\Psif(a)$ is convex and non-increasing.

\section{A Miss Floor for Event-Independent Schedules}\label{sec:floor}
\begin{lemma}[Gap decomposition; restated; proof in the supplement]\label{lem:gap}
Let a deterministic schedule have $n\ge1$ activations at cyclic positions $\tau_1<\dots<\tau_n$ with gaps $g_i=\tau_{i+1}-\tau_i$. Under A1, $1-\mathrm{miss}(x)=T^{-1}\sum_{i=1}^n\phi(g_i)$.
\end{lemma}
\noindent This is the discrete Palm--Khinchin form~\cite{daley2003pp1}, the per-gap cost of patrolling~\cite{lin2013opre}, and for one gap value the periodic quality-of-monitoring formula~\cite[Thm.~4.2]{yau2010tosn}.

\begin{theorem}[Miss floor]\label{thm:floor}
Every event-independent schedule with $\E[a(x)]=a$ has $\mathrm{miss}\ge\eps^*(a)$~\cite[Thms.~4--5]{deep3_jiot2026}.
\end{theorem}
\begin{proof}
For every $d$, $\sum_i\min(g_i,d)\le\min(\sum_ig_i,nd)=\min(T,nd)$. Taking expectations over $D$ and using Lemma~\ref{lem:gap}, $1-\mathrm{miss}\le\E_D[\min(1,(n/T)D)]=\Psif(n/T)$. For a random event-independent schedule, condition on $x$ and apply Jensen's inequality to the concave $\Psif$. The slack is exact: $\mathrm{miss}(x)-\eps^*(a)=T^{-1}\E_D[\min(T,nD)-\sum_i\min(g_i,D)]\ge0$.
\end{proof}
\begin{lemma}[When the floor is tight]\label{lem:tight}
For deterministic $x$, equality holds in Theorem~\ref{thm:floor} iff for every $d$ with $P(D=d)>0$ either all gaps are $\le d$ or all gaps are $\ge d$. Hence equality holds for every $F_D$ supported on $\{1,\dots,T\}$ iff all gaps lie in $\{g,g+1\}$ for some integer $g$. Interleaved schedules with gaps $\lfloor1/a\rfloor$ and $\lceil1/a\rceil$ ($aT\in\mathbb{N}$), and periodic schedules with period $K$ dividing $T$, attain the floor.
\end{lemma}
\begin{proof}
$\sum_i\min(g_i,d)=\min(T,nd)$ iff all $g_i\le d$ (both sides equal $T$) or all $g_i\ge d$ (both sides equal $nd$); if $g_i<d<g_j$ for some $i,j$, the left side is below both. Gaps in $\{g,g+1\}$ leave no integer strictly between two gaps. Conversely, two gaps that differ by at least 2 admit an integer $d\le T$ strictly between them, and a point mass at $d$ breaks equality.
\end{proof}
\noindent The floor is computed from an empirical $F_D$ without knowing the schedule, and Lemma~\ref{lem:tight} shows that it is exact for the schedules a designer would use. Its sufficiency direction is an instance of the optimality of regular (balanced) sequences~\cite{hajek1985moor,altman2000moor}; the converse for every duration law is elementary. For one particular $F_D$ with bounded support, schedules whose gaps all exceed the longest duration also attain it, which is why the equivalence is stated for every duration law. On our corpus $P(D=1)=\pDone$ and \pDltEight{} of events are shorter than 8 frames, so $\eps^*(1/8)=\epsKeight$, although periodic operation with $K=8$ catches every event with $D\ge8$.

\begin{corollary}[Spread beats burst; proof in the supplement]\label{cor:jensen}
Let a schedule be interleaved at rate $a_j$ within each block $j$ of length $n_j$, with $a_jn_j\in\mathbb{N}$ and $\sum_jn_ja_j=aT$. Then $1-\mathrm{miss}\le\Psif(a)$, and $1-\mathrm{miss}=\sum_j(n_j/T)\Psif(a_j)+R$ with $|R|\le J\,\E[D]/T$ for $J$ blocks.
\end{corollary}
\noindent The Jensen gap $\Psif(a)-\sum_j(n_j/T)\Psif(a_j)\ge0$ is the cost of a non-constant rate, up to $J\,\E[D]/T$. It is the principle of~\cite[Thm.~1]{kar2006ton} applied to a harvest-driven rate profile and to events with duration; our earlier article states it qualitatively~\cite{deep3_jiot2026}. Harvest-following, day-only operation and outages are special cases.

\section{Energy Floor and Storage Sizing}\label{sec:energy}
\subsection{The energy-induced floor}\label{sec:efloor}
\begin{lemma}[Storage never empties; restated]\label{lem:moser}
Starting full, a schedule never enters outage iff $\sum_{u=s+1}^t(e_u-h_u)\le C$ for all $s<t$~\cite[Thms.~4,5]{moser2007rts}; leakage enters as in~\cite[Eq.~8]{kansal2007}. The minimal capacity is the maximum drawdown of the cumulative net income.
\end{lemma}
\begin{proof}
Without the lower clip, the recursion with $B_0=C$ has the closed form $B_t=\min_{0\le s\le t}(C+H(s,t]-E(s,t])$, a reflected walk. Every term is non-negative iff the window condition holds. On a cyclic horizon in steady state the minimal $C$ is the maximum drawdown over two concatenated cycles; it is infinite when the cycle's net income is negative.
\end{proof}

Let $\nu_t$ be the least storage at the end of slot $t$ from which idle-only operation survives to the horizon: $\nu_T=0$, $\nu_{t-1}=[\nu_t+1+\lt-h_t]^+$. If $\nu_t>C$ for some $t$, no schedule avoids an outage.

\begin{lemma}[Largest activation budget]\label{lem:greedy}
Let each slot spend activation energy $y_t\in[0,r-1]$ (fluid relaxation). Starting full, the greedy schedule $y_t=\min\{r-1,\,B_{t-1}+h_t-1-\lt-\nu_t\}^+$ maximises $\sum_ty_t$ over all schedules that never enter outage.
\end{lemma}
\begin{proof}
Let $P$ be any feasible schedule. By induction, $B^G_t\le B^P_t$ before the cap: either the greedy stopped at $\nu_t\le B^P_t$, or it spent the maximum while starting from no more storage; $\min(C,\cdot)$ preserves the order. Hence the greedy never spills more than $P$. With $Z_t$ the cumulative spill plus $B_t$, $Z_t=Z_{t-1}+h_t-1-\lt-y_t$ gives $Z^G_T\le Z^P_T$, so $\sum_ty^G_t\ge\sum_ty^P_t$.
\end{proof}
\noindent Every binary schedule that never enters outage is a feasible point of the relaxation, so the relaxed maximum bounds the mean activation of every binary one. Let $A(C)$ be that maximum divided by $T(r-1)$. If in addition the storage at the end of the cycle is at least its start ($B_T\ge B_0$, cyclic steady state), summing the dynamics gives $a\le a_{\mathrm{bal}}=(\bar h-1-\lt)/(r-1)$; let $A_{\mathrm{cyc}}(C)=\min(A(C),a_{\mathrm{bal}})$.

\begin{theorem}[Energy-induced miss floor]\label{thm:energy}
Every event-independent schedule that never enters outage on $(h,C)$ and is in cyclic steady state ($B_T\ge B_0$) has $\mathrm{miss}\ge\eps^*(A_{\mathrm{cyc}}(C))$, even with offline knowledge of $h$. Without the end condition (start at most full, single pass) the bound is $\mathrm{miss}\ge\eps^*(A(C))$.
\end{theorem}
\begin{proof}
Its mean activation is at most $A(C)$ (Lemma~\ref{lem:greedy}; $A$ is non-decreasing in $B_0\le C$), and at most $a_{\mathrm{bal}}$ when $B_T\ge B_0$. $\eps^*$ is non-increasing, and Theorem~\ref{thm:floor} applies.
\end{proof}
\noindent The end condition matters: a start-full single pass can reach $A(C)>a_{\mathrm{bal}}$ when $C$ is large; the simulator repeats the cycle until $B_T=B_0$. For point events ($D=1$), $\eps^*(a)=1-a$ and Theorem~\ref{thm:energy} reduces to capture $\le A_{\mathrm{cyc}}(C)$; for Poisson events, Ren and Cheng~\cite{ren2015p2p} show that activating only when the capacitor is full is optimal. The new content is the spacing term for events with duration. Theorem~\ref{thm:energy} is the first finite-capacity miss floor that holds for every event-independent schedule for events with duration, to our knowledge (Table~S13; Threats (xii)).

\subsection{Outages}\label{sec:blackout}
\begin{proposition}[Blackout]\label{prop:blackout}
An event-independent schedule with outage fraction $\pi$, $N$ outage episodes and on-time activation $a_{\mathrm{on}}$ has $\mathrm{miss}\ge\pi+(1-\pi)\eps^*(a_{\mathrm{on}})-N\E[D]/T$.
\end{proposition}
\begin{proof}
An event that starts in an outage is missed unless it reaches the next on-period, i.e.\ unless its onset lies within $D-1$ slots before an off-to-on boundary; this has probability at most $\E[D]/T$ per boundary. Concatenate the on-slots into one cyclic horizon. An event that starts on an on-slot covers a subset of the concatenated on-slots that its concatenated copy covers, so it is caught only if the copy is caught. Its onset is uniform over on-slots, and Theorem~\ref{thm:floor} bounds these events by $\eps^*(a_{\mathrm{on}})$.
\end{proof}
\noindent At trace scale outages last hours and $\E[D]$ is a few seconds, so the correction is negligible. The term $\pi$ is the loss-of-power-supply probability of storage sizing~\cite{borowy1996lpsp}; only the second term is new. By convexity, $\pi+(1-\pi)\eps^*(a_{\mathrm{on}})\ge\eps^*((1-\pi)a_{\mathrm{on}})$, so an outage costs at least the Jensen gap of Corollary~\ref{cor:jensen}.

\subsection{Sizing for a miss target}\label{sec:sandwich}
\begin{definition}[Sizing quantities]\label{def:sizing}
Let $\eps\in(\eps^*(a_{\mathrm{bal}}),1)$ and $a_\eps=\min\{a:\eps^*(a)\le\eps\}$. $C_{\min}(\eps)$ is the least $C$ at which some event-independent schedule that never enters outage and is in cyclic steady state reaches $\mathrm{miss}\le\eps$. $C_{\mathrm{LB}}(\eps)=\inf\{C:A_{\mathrm{cyc}}(C)\ge a_\eps\}$. $C_{\mathrm{per}}(\eps)$ is the capacity of the interleaved schedule at $\lceil a_\eps T\rceil/T$ (Lemma~\ref{lem:moser}). $C_{\mathrm{TS}}(\eps)$ is the least $C$ at which the realised taut-string slot schedule (below) never enters outage and has exact miss $\le\eps$ by Lemma~\ref{lem:gap}.
\end{definition}
\noindent The taut string comes from the energy-tunnel construction of offline-optimal transmission~\cite{tutuncuoglu2012twc,yang2012tcom}. Write $Y_j$ for the cumulative activation energy up to sub-step $j$ and $S_j$ for the cumulative idle-only surplus. Starting full, the schedules that neither spill nor enter outage are the non-decreasing paths in the tube $S_j\le Y_j\le C+S_j$.
\begin{lemma}[Taut string; variant, proof in the supplement]\label{lem:taut}
In a tube $L\le Y\le U$ with fixed end points, let $\tilde L_j=\max_{i\le j}L_i$ and $\tilde U_j=\min_{i\ge j}U_i$. The taut string $Y^*$ in $[\tilde L,\tilde U]$ satisfies $\sum_j\Psif(\Delta Y^*_j)\ge\sum_j\Psif(\Delta Y_j)$ for every feasible path $Y$ and every concave $\Psif$. It is non-decreasing, and it respects a rate cap whenever some feasible path does.
\end{lemma}
\noindent Neither cited statement covers the monotone hull and the rate cap, so the supplement proves the lemma directly (contact property of $Y^*$ and summation by parts). It also checks it on \LfourN{} trace cells (\LfourOK).

\begin{theorem}[Sizing sandwich]\label{thm:sizing}
For $\eps\in(\eps^*(a_{\mathrm{bal}}),1)$:
(a) $C_{\mathrm{LB}}(\eps)\le C_{\min}(\eps)$;
(b) $C_{\min}(\eps)\le\min\{C_{\mathrm{TS}}(\eps),C_{\mathrm{per}}(\eps)\}$.
\end{theorem}
\begin{proof}
(a) A schedule with miss $\le\eps$ has $a\ge a_\eps$ by Theorem~\ref{thm:floor} and $a\le A_{\mathrm{cyc}}(C)$ by Theorem~\ref{thm:energy}; $A_{\mathrm{cyc}}$ is non-decreasing in $C$.
(b) The realised taut string at any $C$ in the defining set, and the interleaved schedule at $C_{\mathrm{per}}$, are in the class with miss $\le\eps$.
\end{proof}
\begin{remark}[Exactness for a block objective]\label{rem:blk}
Replace the miss by the block objective $1-\sum_j(n/T)\Psif(a_j)$ of Corollary~\ref{cor:jensen}, one block per sub-step, and write $C^{\mathrm{blk}}_{\min}(\eps)$ and $C^{\mathrm{blk}}_{\mathrm{TS}}(\eps)$ for the least capacity and the taut-string capacity under that objective. Let $t^*$ end the sub-step at which the cumulative idle-only surplus is largest. If, at $C^{\mathrm{blk}}_{\min}(\eps)$, some optimal schedule is full at $t^*$ and the non-spilling schedule that dominates it respects the rate cap $a_j\le1$, then $C^{\mathrm{blk}}_{\min}(\eps)=C^{\mathrm{blk}}_{\mathrm{TS}}(\eps)$. Under the cap, a spilling schedule is dominated by a non-spilling one that consumes the would-be spill. A non-spilling cyclic schedule that is full at $t^*$ is a path in the anchored tube, where Lemma~\ref{lem:taut} applies. The hypothesis refers to an optimal schedule and was not checked on the traces. The blocks of $n=600$ slots are also shorter than the longest event (\Dmax{} slots), so the remainder of Corollary~\ref{cor:jensen} is not small. The evaluation therefore uses $C_{\mathrm{TS}}$ with the exact miss; on the traces $C_{\mathrm{TS}}/C^{\mathrm{blk}}_{\mathrm{TS}}$ has median \AiiTSratioMed{} and maximum \AiiTSratioMax.
\end{remark}
\begin{remark}[Policy class]\label{rem:class}
Theorems~\ref{thm:energy} and~\ref{thm:sizing} cover schedules that never enter outage. A schedule that deliberately empties storage saves idle energy during the outage, is still event-independent, and can need less storage than $C_{\mathrm{LB}}$; only Proposition~\ref{prop:blackout} bounds it.
\end{remark}
\noindent Periodic operation attains the least miss at a given realised rate among schedules that never enter outage (Theorem~\ref{thm:floor}, Lemma~\ref{lem:tight}), yet at a given miss it does not need the least storage. The taut string keeps the rate as flat as the energy tube allows: it runs faster while the sun shines and slower before a deficit. Under amendment A2, $C_{\mathrm{LB}}$ is reported as the lower bracket of its bisection and $C_{\mathrm{TS}}$ as a capacity verified feasible, with relative tolerance 0.01. Each bracket lies within that tolerance of its limit, so the convention moves the ratio $C_{\mathrm{LB}}/C_{\mathrm{TS}}$ by at most about two tolerances in either direction.

\section{Trace Evaluation}\label{sec:eval}
\subsection{Setup and pre-registration}\label{sec:setup}
\emph{Harvest.} Satellite global horizontal irradiance for 2025 at native resolution (10--30\,min)~\cite{zippenfeld2023openmeteo}: Ho Chi Minh City (Himawari-9), Madrid and Uppsala (MSG and SARAH-3~\cite{pfeifroth2023sarah3}), four seasons each. Indoor irradiance from EnHANTS setups B--F~\cite{enhants2011dataset}. This gives \nCellsSS{} site--season cells (\nCellHCMC{} + \nCellMadS{} + \nCellUppS{} + \nCellIn). The second satellite product is a sensitivity (Table~S8), and every run also covers it (\nTraceCells{} series--season cells). Traces are held on 1-min sub-steps (zero-order hold; linear interpolation changes $C^*$ by at most \interpMax). The panel is sized so that the mean harvest is $(1+m)$ times the demand of periodic $K=8$, $m\in\{0.25,1\}$.

\emph{Events.} The running example is a camera node on solar or indoor light~\cite{desai2022camaroptera,mayer2024iotj}. The duration law $F_D$ is rebuilt from detector outputs~\cite{ultralytics_yolo} on \nVid{} public videos (CDnet~2014~\cite{wang2014cdnet}, LASIESTA~\cite{cuevas2016lasiesta}, BMC~\cite{vacavant2013bmc}). An event is a maximal run of frames with at least one detection; there are \nEv{} events with $\E[D]=\ED$ frames. LASIESTA frames are ordered by the number in their file name. An earlier build was in error: it cut LASIESTA events on the text order of the names, which gave \VoneNev{} events (\VoneLas{} LASIESTA events instead of \VtwoLas). Reference~\cite{deep3_jiot2026} reports its statistics on that mis-ordered build: after correction its printed floors rise ($\eps^*(1/5)$ from \FloorFiveOld{} to \FloorFiveNew), its spacing conclusions stand (Lemma~\ref{lem:tight}), and its gate advantage stays marginal (Section~S1). The correction is amendment A1, frozen before any corrected result existed; the supplement compares both versions. The correction changed one verdict: H3a passed on the earlier corpus and fails in magnitude on the corrected one (Table~S1).

\emph{Grid and protocol.} The grid crosses $r\in\{1.1,\dots,300\}$ (8 values), $\lt\in\{0,0.01,0.04,0.1\}$, $m$, and 20 capacities per combination, with eleven policies (six periodic with $K\in\{2,\dots,64\}$, throttling, harvest-following, a Kansal-style EWMA duty cycle, taut string, greedy). Odd months and half of the videos (seeded) form TRAIN; even months and the other videos form TEST. Predictions with pass criteria were frozen (SHA-256 and a git tag) after TRAIN and before TEST. Every number below is from TEST. Misses are measured by replaying $10^5$ real events per run; bootstrap intervals resample videos ($B=1000$) or 7-day blocks of the trace. Sign tests use Holm's correction at family-wise $\alpha=0.05$. The \nCombos{} combinations of the sizing-rule comparison are \nCells{} series--season cells times \nRgrid{} values of $r$, \nLgrid{} of $\lt$ and \nMgrid{} of $m$. The sizing ratios do not depend on $r$ (Table~S12), so the effective sample is \nCellsSS{} site--season cells: \nCellHCMC{} in Ho Chi Minh City, \nCellMadS{} in Madrid, \nCellUppS{} in Uppsala and \nCellIn{} indoor setups. Table~\ref{tab:prereg} lists every prediction, its criterion and its outcome on the corrected corpus. H1, H2a and H6 are consistency checks: they compare a proved statement or an identity with a simulator of the same model.

\providecommand{\preregCritW}{0.50\textwidth}\providecommand{\preregObsW}{0.335\textwidth}
\input{prereg_table.tex}

\subsection{The floors hold and are tight}\label{sec:holds}
Across \nTone{} policy--capacity runs, no replayed miss lies below the floor of Theorem~\ref{thm:floor} by more than a Bonferroni margin ($z=\zBonf$). The same holds for the \nPtwo{} blackout runs (Proposition~\ref{prop:blackout}) and the \nTtwo{} no-outage runs (Theorem~\ref{thm:energy}); the latter also have no analytic violation above $10^{-6}$. In Fig.~\ref{fig:floor}, periodic and taut-string runs lie on the floor within replay noise; analytically, periodic schedules with $K$ dividing $T$ attain it to within $10^{-12}$ (supplement). Points above the floor have a variable rate or outages.

\begin{figure}[!t]\centering
\includegraphics[width=\columnwidth]{figures/F2_miss_vs_floor_test.pdf}
\caption{Replayed miss against the applicable floor (Theorem~\ref{thm:floor}, \ref{thm:energy} or Proposition~\ref{prop:blackout}) for a random sample of \FtwoN{} TEST runs. \Ftwocaption}\label{fig:floor}
\end{figure}

\subsection{Spreading versus following the harvest}\label{sec:spread}
At equal mean rate $a=1/8$, harvest-following needs \dCMinS--\dCMaxS{} less storage than periodic operation but misses \dmissMinS--\dmissMaxS{} more events, significant in all \nCellsSS{} cells after Holm over \nHolmS{} tests (largest $p=\pMaxBootS$). The family is \nHolmHthree{} H3c tests on these cells plus \nHolmHfourd{} H4d tests on the primary cells among the registered six. The sign held in every combination, but H3a failed in magnitude: the medians at $K=8$ are \HthreeaLo--\HthreeaHi{} against a registered band of $[0.18,0.32]$. At equal miss target the comparison reverses. With the exact brackets of A2, the taut string needs only \AiiTsMedLo--\AiiTsMedHi{} of the periodic capacity (medians per site group, $\eps\in\{0.5,0.6\}$), below the registered band of $[0.2,0.7]$ for Uppsala and indoor (H4a, Table~\ref{tab:prereg}). The registered values with the block objective are \tsMedLo--\tsMedHi{} over \nEps{} cells (Table~\ref{tab:prereg}). Across \AiiNfinite{} finite cells no lower bracket exceeds the verified feasible capacity (largest ratio \AiiMaxLBTS); the exact and block misses of the taut string differ by a median of \AiiAbsMed{} (maximum \AiiAbsMax). The lower bound is within \lbtsGapPct{} of $C_{\mathrm{TS}}$ for $\eps\ge0.6$ and loose at $\eps=0.45$ (\lbtsLoMin--\lbtsLoMax). H4a--H4d keep their registered verdicts; their A2 column is in the supplement.

\begin{figure}[!t]\centering
\includegraphics[width=\columnwidth]{figures/F3_frontier_spread_vs_hf_test_a2_v3.pdf}
\caption{Left: harvest-following against periodic at $a=1/8$ ($r=15$, $\lt=0.04$, $m=0.25$; one point per cell). Right: taut-string capacity relative to periodic at equal miss target (median per site, exact brackets of A2; the registered block-objective curves, dashed, coincide).}\label{fig:frontier}
\end{figure}

\subsection{Sizing}\label{sec:sizing}
Table~\ref{tab:sizing} gives the storage needed by the reference node ($r=14.3$, $e_{\mathrm{idle}}=0.248$\,mJ, $\lt=0.036$, $K=8$; \Pidle\,mW idle, \Pact\,mW during an activation of \Tact\,ms) over the full trace. At $m=0.25$ the implied panel is \WpHCMCLo--\WpHCMCHi\,mW$_p$ in Ho Chi Minh City, \WpMadLo--\WpMadHi\,mW$_p$ in Madrid and \WpUppLo--\WpUppHi\,mW$_p$ in Uppsala. It follows from $W_p=P_{\mathrm{mean}}\cdot1000/\overline{\mathrm{GHI}}$ for a horizontal panel of constant efficiency, without temperature, MPPT or regulator loss. The baselines are average-day sizing, Buchli's daily-resolution deficit~\cite{buchli2014sensys}, Kansal's condition with the mean harvest rate~\cite{kansal2007}, and a 3-day autonomy rule. Across the full grid, average-day sizing underestimates $C^*$ by a median factor of \avgHCMC{} in Ho Chi Minh City (\nCellHCMC{} cells). The factor is \avgMadS{} in Madrid (\nCellMadS{} cells; \avgMadSec{} with the second satellite product) and \avgUppS{} in Uppsala (\nCellUppS{} cells; \avgUppSec{} with the second product). The outdoor maximum is large (\MaxOutPrim{} with the primary product, \avgMaxOut{} with either). Indoors the median is \avgIn{} over \nCellIn{} setups (per-setup medians \avgInSetupLo--\avgInSetupHi) and the maximum \avgMaxIn. These medians hide a wide spread between cells (amendment A3, post hoc; Table~S15). The per-cell medians run from \CellHCMCLo{} to \CellHCMCHi{} in Ho Chi Minh City, \CellMadLo{} to \CellMadHi{} in Madrid and \CellUppLo{} to \CellUppHi{} in Uppsala. The full-trace ratios in Table~\ref{tab:sizing} are larger, \FullRatioLo--\FullRatioHi{} in Ho Chi Minh City, because the full-trace $C^*$ is set by the season's worst spell. Daily resolution hides the night in the tropics (Buchli median \buHCMC{} in Ho Chi Minh City, on the TEST grid; Table~\ref{tab:sizing}'s full-trace ratios are \BuFullLo--\BuFullHi; \buLoS--\buHiS{} elsewhere). Buchli's ratio is undefined in the \nBuDrop{} combinations of Madrid in summer with $m=1$, where every day's harvest covers that day's demand, so H2d counts \nBuN{} combinations. The 3-day rule always oversizes in Ho Chi Minh City (largest ratio \threeHCMCmax) but undersizes Uppsala by up to \MaxThreeUppPrim{} times (\threeUppMax{} with either product). With the mean harvest rate, Kansal's condition oversizes by \kmeanInvLo--\kmeanInvHi{} times. The table is descriptive; its pre-registered one-month version failed in magnitude (H7), because $C^*$ of a single month depends on that month's worst spell. For the same reason Table~\ref{tab:sizing} shows the second satellite product and, as a post-hoc check under amendment A3, $C^*$ without the 7-day window of least net income. Over the \nCellsSS{} cells that ratio has median \WorstMedS{} (range \WorstMinS--\WorstMaxS). In \WorstOneS{} cells $C^*$ is unchanged because that week is not the spell that sets it (Table~S14). The period is the cheapest lever in winter. Doubling $K$ from 8 to 16 cuts the winter need by a median factor of \ckWinter{} at the European sites and \ckHCMC{} in Ho Chi Minh City. The miss rises from \epsKeight{} to $\eps^*(1/16)=\epsKsixteen$.

\emph{Report energy (amendment A4, post hoc).} The registered $r=14.3$ holds no report energy, because its source measures only sensing, so $r=14.3+E_{\mathrm{rep}}/e_{\mathrm{idle}}$. The models of~\cite{vanleemput2023iotj} give, for one 8-byte uplink report, $E_{\mathrm{rep}}=\ErepBle$\,mJ on BLE Mesh, \ErepSix\,mJ on 6TiSCH and \ErepLora\,mJ on LoRaWAN Class~A at SF7, so $r$ is \rBle, \rSix{} and \rLora. The panel is rescaled with the demand, so $C^*$ grows by one factor in every cell: \kBle, \kSix{} and \kLora{} (Table~S17). With a rescaled panel the result is a multiplier and no new trace effect. No lower bracket exceeds the verified capacity in the \AivNcases{} cases run (largest ratio \AivMaxLB). LoRaWAN is an energy-only figure: factor \kLora, not schedulable at $K=8$, because its Class~A task of \TaskLora\,s exceeds the activation period of \ActPeriod\,s.

\emph{Fixed panel (amendment A5, post hoc).} A planner who keeps the panel pays more, by an amount the trace sets. The registered panel has a margin of \PanelMarginQ\,mJ per activation, so BLE and 6TiSCH keep the energy balance in every cell. On that panel 6TiSCH needs \AvSixMin--\AvSixMax{} times the storage of the node without reports (median \AvSixMed{} over \nCellsSS{} cells), against \kSix{} with a rescaled panel. The factor is smallest for indoor setup~F and largest in Ho Chi Minh City in June--August; with the second satellite product it is \AvSixSecMin--\AvSixSecMax. BLE needs at most \AvBleMax{} times. No lower bracket exceeds the verified capacity in the \AvNcases{} fixed-panel cases (largest ratio \AvMaxLB).

\input{sizing_table.tex}

\subsection{Throttling, blackout and spending the spill}\label{sec:throttle}
Throttling lowers the activation rate when storage is low instead of letting the node black out. In the \nThr{} TEST cases where periodic $K=8$ has outages, even the best of five throttling variants per case (chosen with hindsight) stays within $\pm0.03$ of blackout in all cases (median \thrMedSigned). The common rule ``drop from $K=8$ to $K=64$ below 20\,\% storage'' misses more than blackout in \naiveSharePct{} of cases (median \naiveMedSigned). Outages last hours, so a throttled node still pays the idle energy and enters the same deficit.

\emph{Exploratory: a predictive controller.} A Kansal-style predictive duty cycle~\cite{kansal2007} misses fewer events than blackout at the same storage (median \ExKblk). Three follow-up predictions registered before a dedicated post-test stage all failed. In particular, against periodic operation set to the controller's measured mean activation, the predicted sign held in only \ExOneShare{} of cases (median \ExOneMed; opposite sign significant in \ExOneWrong{} of \ExOneN{} series--season cells). This comparison (E1) matched nominal rates. Because periodic operation enters outage in every case of E1 while the controller spends the surplus, the two run at different realised rates and energies, so E1 does not test Theorem~\ref{thm:floor}. A per-case check, run post hoc, confirms this reading. The controller is better by more than 0.005 in \KcNwin{} of \KcNcases{} cases; there the realised periodic rate is \KcRateRatio{} times lower and its floor \KcGapMed{} higher (medians). The controller has more outage than periodic operation in \KcKmoreOutPct{} of cases, but runs faster while on. It never misses less than $\eps^*$ of its own realised rate (\KcKoneFail{} of \KcNall{} cases below by more than 3 standard errors). Where periodic operation has no outage (\KcNnoOut{} cases, outside E1), the controller misses more in \KcNoOutLosePct{} of them, as Theorem~\ref{thm:floor} predicts. The comparison at equal storage is E3. A constant rate sized on the TRAIN months to never black out did worse than both the controller and blackout (median \ExThreeBlk{} against blackout; better than both in \ExThreeBoth{} of cases). It had outages in \ExThreeOutZero{} of the TEST cases, and still in \ExThreeOutTen{} with a 10\,\% rate margin.

\section{A Sizing Procedure for Deployments}\label{sec:procedure}
The results combine into a procedure that a fleet planner runs per site before deployment, to choose a supercapacitor and a duty cycle, or to check a miss claim of an event-independent, never-outaged controller. Every step is a closed form or a maximum drawdown on a measured trace.
\begin{enumerate}
\item Estimate $F_D$ from a sample of event durations and tabulate $\eps^*(a)$; the miss target fixes the activation rate $a_\eps$.
\item The energy balance bounds the reachable floor by $\eps^*(a_{\mathrm{bal}})$; a target below it needs a larger panel rather than more storage.
\item Size storage with $C_{\mathrm{per}}(\eps)$ (Lemma~\ref{lem:moser}) for a constant-rate node, or with $C_{\mathrm{TS}}(\eps)$ for a node that knows its harvest. On the measured trace, Theorem~\ref{thm:sizing} brackets every event-independent controller that never enters outage; a claim of such a controller to need less than $C_{\mathrm{LB}}(\eps)$ on that trace is inconsistent with the model.
\item Size on the worst season of several years of data, and keep a margin. $C^*$ is set by the single worst spell, and the bound holds only for the trace it was computed on. On our traces a 10\,\% rate margin was not enough: a constant rate set on the TRAIN months still had outages in \ExThreeOutTen{} of TEST cases (Section~\ref{sec:throttle}). A storage margin behaves similarly (amendment A3, post hoc; Table~S16). With the panel fixed on the TRAIN months, the capacity needed on the TEST months is a median of \KMedS{} times the TRAIN capacity. It exceeds twice that capacity in \KGtTwoS{} of \nCellsSS{} cells, and in \KInfS{} further cell (with the second satellite product; Table~S16) no capacity suffices, which calls for a larger panel. On 2020--2025 in Ho Chi Minh City the seasonal $C^*$ varied by a factor of \AivbLo--\AivbHi{} (amendment A4b, post hoc; Table~S18). A season of 2025 needed up to \AivbOverFive{} times less storage than the same season of the worst year, so size on the worst year. The procedure therefore delivers a number without a proved margin; the margin must come from multi-year data.
\item Do not rely on threshold throttling to soften outages; add storage or panel. Spending the spill with a predictive duty cycle helped in the exploratory analysis of Section~\ref{sec:throttle}; this step is exploratory; no bound covers it.
\item Before adding a content-based gate, check its cue cost against the boundary of Proposition~\ref{prop:gate}.
\end{enumerate}
\noindent\emph{Worked example.} Take the reference node in Ho Chi Minh City in September--November, reporting over 6TiSCH ($r=\rSix$, $K=8$, $m=0.25$). Its mean harvested power is \ExPmean\,mW, which needs a panel of \ExWp\,mW$_p$. The third step gives $C^*=\ExCkJ$\,kJ (\ExCkJFixed\,kJ with the panel held fixed), about \ExCaps{} supercapacitors of 10\,F at 3\,V. The fourth step applies: for the registered $r$ the worst season of 2020--2025 needed \AivbWorstkJ\,kJ, against \AivbFivekJ\,kJ for the worst season of 2025 alone.

\begin{proposition}[Gate boundary]\label{prop:gate}
Let periodic operation at rate $a_{\mathrm{per}}$ be exactly sustained at capacity $C^*_{\mathrm{per}}>0$. A gate that activates at long-run rate $\rho\,a_{\mathrm{per}}$, spread over time, and pays $c$ per slot empties storage at $C^*_{\mathrm{per}}$ iff $c>c^*(\rho)=(1-\rho)a_{\mathrm{per}}(r-1)$.
\end{proposition}
\begin{proof}
The gate changes the demand per slot by $c+(\rho-1)a_{\mathrm{per}}(r-1)$, and $C^*$ is strictly increasing in a constant demand when $C^*_{\mathrm{per}}>0$ (Lemma~\ref{lem:moser}).
\end{proof}
\noindent The boundary does not depend on the site. An ideal presence gate on our corpus would be active on \inEventFrac{} of frames ($\rho\approx\rhoPresence$ against $K=8$), so it is infeasible at $C^*_{\mathrm{per}}$ for every $c\ge0$. The simulated map agrees with the boundary in \nGateAgree{} of \nGate{} cells (supplement). A bursty gate is outside the proposition in both directions.

\section{Threats to Validity}\label{sec:threats}
(i) Onsets are assumed uniform (A1). Real event rates are diurnal, which weights $\Psif$ by time of day and breaks the site independence of $\eps^*$; the corpus has no wall-clock timestamps, so we could not test this. (ii) An event is a maximal run of detections, so $F_D$ is the run-length law of the detector on surveillance video rather than of the scene. For such events the assumption that an invoked detector catches the event holds by construction. $F_D$ is from a different detector than the reference node's. Only $F_D$ is used, and the floor is computable for any other duration law. (iii) Theorems~\ref{thm:energy} and~\ref{thm:sizing} do not cover schedules that deliberately empty storage (Remark~\ref{rem:class}); sizing for that class is open. (iv) Remark~\ref{rem:blk} rests on an anchoring hypothesis that we did not check; the gap between the block objective and the exact miss is measured on the traces and has no proved bound. (v) The restart threshold $\theta$ enters outage lengths but no theorem; we did not vary it. (vi) Five analysis choices were made after results were seen and are disclosed as amendments. The frame-order check of A1 used a graded overlap measure, chosen after the registered hit-rate clause was not met at $\pm1$ frame; A2 replaced the reporting of the sandwich brackets. A3 added the worst-spell, per-cell and margin sensitivities, and A4 the report-energy sensitivity, the count on \nCellsSS{} cells and the inter-annual range, and A5 the fixed-panel variant; they are descriptive and change no verdict. (vii) In Uppsala, $C^*$ from MSG and SARAH-3 differs by up to \uppProdDiffMax{} (median \uppProdDiffMed), so we report both. (viii) $C^*$ is a maximum drawdown, and block-bootstrap intervals break the worst spell and can exclude the point estimate; see \blockNote. (ix) The taut string needs the harvest in advance and serves as an offline bound; no hardware was measured. (x) The report energy of one node is inside the model through $r$, as a worst case of one report per activation; network-level effects (collisions, retransmissions) are outside these bounds. (xi) One calendar year of satellite data for two sites; inter-annual range measured for one site only, with the panel rescaled each year. (xii) Positioning against~\cite{ren2014tpds},~\cite{ren2015p2p},~\cite{ballestrin2026cl} rests on their abstracts; full texts were not consulted.

\section{Conclusion}\label{sec:conclusion}
Two closed-form quantities, the miss floor $\eps^*(a)$ and the energy-limited activation $A_{\mathrm{cyc}}(C)$, bound how often an energy-harvesting sensing node that never enters outage can catch events, and how much storage it needs. On held-out traces the floors are never violated and evenly spaced schedules attain them. Common sizing rules are off by large factors: average-day sizing by a median of up to \avgIn{} times indoors, Kansal's mean-rate condition by \kmeanInvLo--\kmeanInvHi{} times the other way. In practice, size from the duration law and the worst spell of several years of measured trace, spread activations evenly, and do not expect threshold throttling to soften outages. At one site the seasonal need varied up to \AivbHi{}-fold over six years. The half-year-to-half-year factor had median \KMedS{} and exceeded two in \KGtTwoS{} of \nCellsSS{} cells, so the procedure gives a number without a proved margin. A report on every activation multiplies the need by a factor set by the radio if the panel grows with the demand. If the panel stays, the factor on 6TiSCH is \AvSixMin--\AvSixMax, depending on the trace.

\section*{Data Availability}
Code (MIT), the pre-registration with its amendments and the derived data (CC BY 4.0) that generate every number, table and figure are at \repoURL{}; the event corpus holds onsets and durations only. The third-party inputs (Section~\ref{sec:setup}) are not redistributed.

\section*{Acknowledgment}
The author used AI assistants by Anthropic. Claude Opus 5.5 was used for drafting the text, code development and literature search, and Claude Fable 5.1 for a second check of the proofs and a review of the manuscript. All results, code and text were checked by the author, who takes full responsibility for them.

\bibliographystyle{IEEEtran}
\bibliography{refs}

\begin{IEEEbiographynophoto}{Dat Lam Quoc}
is a researcher with the FPT School of Business and Technology (FSB), FPT University, Ho Chi Minh City, Viet Nam. Research interests include energy-aware sensing and edge inference. ORCID 0009-0004-5432-9343.
\end{IEEEbiographynophoto}
